\newpage
% =================================================
\section{Tecnologia e Relatórios (Vercel / Firebase)}
% =================================================

\subsection{Para relatórios}
Após avaliação das restrições do ambiente \textit{serverless} (Vercel e Firebase Functions), a biblioteca escolhida para a geração de relatórios no backend é o \textbf{\textit{pdfkit}} em conjunto com o plugin \textbf{\textit{pdfkit-table}}.

\textbf{Justificativa da escolha:}
\begin{itemize}
    \item \textbf{Baixo consumo de memória:} O \textit{pdfkit} desenha o PDF nativamente, respeitando os limites de memória da Vercel/Firebase.
    \item \textbf{Tabelas dinâmicas:} O \textit{pdfkit-table} resolve a paginação automática de tabelas longas.
    \item \textbf{Fluxo de Armazenamento:} A Cloud Function gera o PDF em um \textit{Buffer} de memória, retorna o PDF em base64 conforme o contrato atual; o cliente cria um Blob temporário para download.
\end{itemize}

\subsubsection{Assinatura Lógica de Relatórios (Hash Canônico)}
Todos os relatórios gerados pelo backend injetam um rodapé contendo um Hash Canônico. Este hash compila os parâmetros da requisição (ID, datas e registros consultados) concatenados com o Timestamp (\texttt{Date.now()}), servindo como identificador de rastreabilidade, sem equivaler a assinatura digital.

\subsubsection{Prevenção de Limites de Leitura Firestore (N+1)}
Resolver referências únicas em lotes para reduzir N+1. O tamanho é definido pela API e pelo volume medido; in e getAll têm contratos distintos, sem limite universal de 200 itens.

\subsubsection{Módulo de Geração de Etiquetas (PDF)}
Para a identificação física dos frascos, o backend também utiliza o \textbf{pdfkit} para desenhar uma matriz (Grid Visual de Offset 3×10) em tamanho A4, com etiquetas de 65,0 mm × 26,5 mm. O código de barras no formato Code 128 é gerado ativamente em memória via biblioteca \textbf{bwip-js} e embutido no buffer do PDF, garantindo alta definição na impressão e dispensa de fontes especiais no client. Para frascos já cadastrados (2ª via), o backend gera folhas individuais A4 contendo a Ficha de Conferência e Rastreabilidade com histórico recente de 10 movimentações e 1 única etiqueta de reposição centralizada no rodapé.

\begin{explicacao}[Leitura dos exemplos legados]
Os listings a seguir preservam exemplos de desenvolvimento, não uma implementação pronta para copiar. O contrato atualizado está nas seções 5.9, 8 e 11: caminhos de especificações são aninhados, datas civis são normalizadas no fuso institucional, relatórios usam base64, consultas e escritas são ordenadas em transação e claims exigem controle de revogação. Q06 utiliza tolerância híbrida, consumo não negativo e ajuste separado quando há ganho de massa tolerado. Não adotar consumo negativo, data de abertura inventada, placeholder de foto ou simples cópia de novo nome sobre o bem. Esses exemplos devem ser adaptados ao contrato antes de implementação.
\end{explicacao}

\subsection{Princípio: nunca confiar no \textit{client-side}}
Todas as \textit{Cloud Functions} desta seção seguem a mesma regra: \textbf{qualquer dado que o servidor consiga derivar ou verificar por conta própria é recalculado ou reconferido no servidor, nunca aceito como o cliente enviou}. Isso vale, em particular, para o papel do usuário (claims assinadas pelo Auth, nunca um papel declarado no payload; conta ativa verificada seletivamente conforme DP-D01), para valores calculados como peso/volume/vencimento/atraso (sempre recomputados a partir dos dados brutos e da hora do servidor), para códigos gerados pelo sistema (Regra 6.2.9), e para o estado de disponibilidade de um recurso concorrente (sempre relido dentro de uma transação antes de decidir).

\subsubsection{Funções de Apoio (Autorização)}

\begin{lstlisting}[language=TypeScript, title={Resolução de Papéis via Custom Claims (zero leituras Firestore)}]
import { onCall, HttpsError } from "firebase-functions/v2/https";
import * as admin from "firebase-admin";

// C1: Papéis lidos diretamente do JWT (Firebase Custom Claims).
// Zero leituras extras no Firestore por chamada de API.
// PREREQUISITO: toda mutação de papel deve chamar atualizarCustomClaims(uid).
function resolverPapeisDoToken(auth: {token: Record<string, unknown>} | undefined): string[] {
  if (!auth) return [];
  const roles = auth.token["roles"];
  if (!Array.isArray(roles)) return [];
  return roles as string[];
}

// Chamada após TODA concessão ou remoção de papel:
//   await atualizarCustomClaims(uid);
// Isso atualiza claims para o próximo token; não invalida automaticamente o atual; o cliente deve renovar com user.getIdToken(true).
export async function atualizarCustomClaims(uid: string): Promise<void> {
  const colecoes = ["Chefe_Geral", "Gestor_Almoxarifado", "Gestor_Bens_Patrimoniais", "Professor", "Aluno", "Bolsista"];
  const leituras = await Promise.all(
    colecoes.map((c) => admin.firestore().collection(c).doc(uid).get())
  );
  const roles = colecoes.filter((_, i) => leituras[i].exists);
  await admin.auth().setCustomUserClaims(uid, { roles });
}

async function validarPermissao(
  request: {auth?: {uid: string; token: Record<string, unknown>}},
  papeisPermitidos: string[],
  requerAtivo = false
): Promise<string[]> {
  if (!request.auth?.uid) throw new HttpsError("unauthenticated", "Usuário não autenticado.");
  if (requerAtivo) {
    const snap = await admin.firestore().collection("Usuarios").doc(request.auth.uid).get();
    if (!snap.exists || snap.data()?.ativo !== true)
      throw new HttpsError("permission-denied", "Conta desativada ou não habilitada.");
  }
  const papeis = resolverPapeisDoToken(request.auth);
  if (!papeisPermitidos.some((p) => papeis.includes(p)))
    throw new HttpsError("permission-denied", "Usuário sem papel autorizado para esta ação.");
  return papeis;
}

async function validarGestorDoAlmoxarifado(uid: string, token: Record<string, unknown>, idAlmoxarifado: string): Promise<void> {
  const papeis = resolverPapeisDoToken({ token });
  if (papeis.includes("Chefe_Geral")) return; // Chefe tem escopo global
  if (!papeis.includes("Gestor_Almoxarifado")) {
    throw new HttpsError("permission-denied", "Usuário não é Gestor de Almoxarifado.");
  }
  // 1 leitura: o vínculo de escopo (não o papel - esse já veio do token)
  const vinculo = await admin.firestore().collection("Gestor_Almoxarifado_x_Almoxarifado")
    .where("id_gestor_almoxarifado", "==", uid)
    .where("id_almoxarifado", "==", idAlmoxarifado)
    .limit(1).get();
  if (vinculo.empty) {
    throw new HttpsError("permission-denied", "Gestor não está designado para este almoxarifado.");
  }
}

// Geração atômica sequencial no Firestore via Documento Singleton (Seção 5.7)
// Executado estritamente dentro da transação de cadastro do frasco
async function obterProximoCodigoFrascoTx(
  tx: FirebaseFirestore.Transaction
): Promise<{ proximoNumero: number; novoCodigo: string }> {
  const contadorRef = admin.firestore().collection("Contador_Codigo_Frasco").doc("singleton");
  const contadorSnap = await tx.get(contadorRef);

  const ultimoCodigo = contadorSnap.data()?.ultimo_codigo_gerado || 0;
  const proximoNumero = ultimoCodigo + 1;

  tx.set(contadorRef, { ultimo_codigo_gerado: proximoNumero }, { merge: true });

  return { proximoNumero, novoCodigo: `LCQUI-${proximoNumero}` };
}
\end{lstlisting}

\begin{regra}[DP-D01: verificação seletiva e risco residual]
Usar \texttt{requerAtivo=true} em retirada, devolução, baixa patrimonial, aprovação/rejeição de requisições e concessão/revogação de papéis. As funções \texttt{concederPapel} e \texttt{revogarPapel} seguem a mesma assinatura. Leituras e mutações sem essa exigência podem continuar confiando em token válido até renovação (aproximadamente uma hora). Verificar ativo não atualiza automaticamente papéis de uma conta ainda ativa. Essa limitação é risco residual aceito da DP-D01, não revogação universal instantânea.
\end{regra}

\subsubsection{Fluxo de Reagentes: Cadastro, Retirada e Devolução}

\begin{lstlisting}[language=TypeScript, title={Cadastro de Frasco Fechado (Regra 6.2.13, item 1)}]
interface CadastroFrascoFechado {
  idResumoReagente: string;
  idEspecificacaoReagente: string;
  idAlmoxarifado: string;
  idLote?: string;
  pesoTotal: number;
  volumeNominal: number;
  validadeFechado?: string;
  validadeDesconhecida?: boolean;
  decisaoSeJaVencido?: "QUARENTENA" | "PENDENTE_DE_DESCARTE" | "DISPONIVEL";
  detalheStatus?: string;
}

export const cadastrarFrascoFechado = onCall(async (request) => {
  const dados = request.data as CadastroFrascoFechado;
  await validarPermissao(request, ["Chefe_Geral", "Gestor_Almoxarifado"]);
  await validarGestorDoAlmoxarifado(request.auth!.uid, request.auth!.token, dados.idAlmoxarifado);

  const especSnap = await admin.firestore().collection("Resumo_Reagente").doc(dados.idResumoReagente).collection("Especificacoes").doc(dados.idEspecificacaoReagente).get();
  if (!especSnap.exists) throw new HttpsError("not-found", "Especificação de reagente não encontrada.");
  const densidade = especSnap.data()!.densidade;
  const resumo = (await admin.firestore().collection("Resumo_Reagente").doc(dados.idResumoReagente).get()).data()!;
  const estadoFisico = resumo.estado_fisico;

  let pesoVazio: number;
  if (estadoFisico === "LIQUIDO") {
    if (!densidade || densidade <= 0) throw new HttpsError("failed-precondition", "Líquido exige densidade positiva.");
    pesoVazio = dados.pesoTotal - (dados.volumeNominal * densidade);
  } else {
    pesoVazio = dados.pesoTotal - dados.volumeNominal; // g
  }

  if (pesoVazio <= 0) {
    throw new HttpsError("invalid-argument", "Peso total incompatível com o volume nominal para esta densidade.");
  }

  const agora = new Date();
  const validadeFechado = dados.validadeFechado ? new Date(dados.validadeFechado) : null;
  const validadeEfetiva = dados.validadeDesconhecida ? null : validadeFechado;
  const venceuNoCadastro = Boolean(validadeEfetiva && validadeEfetiva <= agora);

  if (venceuNoCadastro && !dados.decisaoSeJaVencido) {
    throw new HttpsError("failed-precondition", "O frasco está vencido no cadastro e exige uma decisão do gestor.");
  }

  if (dados.decisaoSeJaVencido === "QUARENTENA" && !dados.detalheStatus) {
    throw new HttpsError("invalid-argument", "A entrada em quarentena exige detalhe_status.");
  }

  // Execução atômica no Firestore: Sequenciamento + Cadastro do Frasco
  return admin.firestore().runTransaction(async (tx) => {
    // 1. Gera o código sequencial estrito LCQUI-N
    const { novoCodigo } = await obterProximoCodigoFrascoTx(tx);

    // 2. Persiste o frasco com o código gerado
    const frascoRef = admin.firestore().collection("Frasco_Reagente").doc();
    tx.set(frascoRef, {
      id_almoxarifado: dados.idAlmoxarifado,
      id_lote: dados.idLote ?? null,
      id_especificacao_reagente: dados.idEspecificacaoReagente, // projeção física
      id_resumo_reagente: dados.idResumoReagente,
      eh_higroscopico: resumo.eh_higroscopico, // snapshot servidor imutável
      codigo_frasco: novoCodigo,
      conteudo_nominal: dados.volumeNominal,
      peso_no_cadastrado: dados.pesoTotal,
      peso_atual: dados.pesoTotal,
      peso_frasco_vazio: pesoVazio,
      medida_usada: 0,
      estado_fisico_frasco: "FECHADO",
      disponibilidade: "DISPONIVEL",
      validade_fechado: validadeEfetiva,
      validade_efetiva: validadeEfetiva,
      validade_desconhecida: dados.validadeDesconhecida ?? false,
      vencido: venceuNoCadastro,
      em_quarentena: venceuNoCadastro && dados.decisaoSeJaVencido === "QUARENTENA",
      uso_vencido_autorizado: venceuNoCadastro && dados.decisaoSeJaVencido === "DISPONIVEL",
      detalhe_status: venceuNoCadastro
        ? dados.detalheStatus ?? (dados.decisaoSeJaVencido === "PENDENTE_DE_DESCARTE" ? "Frasco vencido aguardando processo institucional de descarte." : null)
        : null,
      cadastrado_em: admin.firestore.FieldValue.serverTimestamp(),
      cadastrado_por: request.auth!.uid,
    });

    return { idFrasco: frascoRef.id, codigoFrasco: novoCodigo, pesoVazioCalculado: pesoVazio, venceuNoCadastro };
  });
});
\end{lstlisting}

\begin{lstlisting}[language=TypeScript, title={cadastrarFrascoAberto com bifurcação SÓLIDO/LÍQUIDO (C5)}]
// C5: modalidade para SÓLIDO usa ESTIMA_MASSA (não ESTIMA_VOLUME)
interface CadastroFrascoAberto {
  idResumoReagente: string;
  idEspecificacaoReagente: string;
  idAlmoxarifado: string;
  idLote?: string;
  // LIQUIDO: "CONHECE_TARA" | "ESTIMA_VOLUME"
  // SOLIDO:  "CONHECE_TARA" | "ESTIMA_MASSA"
  modalidade: "CONHECE_TARA" | "ESTIMA_VOLUME" | "ESTIMA_MASSA";
  pesoTotalBalanca: number;            // sempre leitura de balança em g
  pesoFrascoVazioInformado?: number;  // Opção A (ambos estados físicos)
  volumeAtualEstimado?: number;        // Opção B LIQUIDO: volume visual em mL
  massaAtualEstimada?: number;         // Opção B SOLIDO: estimativa em g
  dataAbertura?: string;
  aberturaHistoricaDesconhecida: boolean;
  validadeAberto?: string;
  validadeDesconhecida?: boolean;
  decisaoSeJaVencido?: "QUARENTENA" | "PENDENTE_DE_DESCARTE" | "DISPONIVEL";
  detalheStatus?: string;
}

export const cadastrarFrascoAberto = onCall(async (request) => {
  const dados = request.data as CadastroFrascoAberto;
  const papeis = await validarPermissao(request, ["Chefe_Geral", "Gestor_Almoxarifado"]);
  await validarGestorDoAlmoxarifado(request.auth!.uid, request.auth!.token, dados.idAlmoxarifado);

  const especSnap = await admin.firestore().collection("Resumo_Reagente").doc(dados.idResumoReagente).collection("Especificacoes").doc(dados.idEspecificacaoReagente).get();
  if (!especSnap.exists) throw new HttpsError("not-found", "Especificação não encontrada.");
  const especData = especSnap.data()!;
  const resumo = (await admin.firestore().collection("Resumo_Reagente").doc(dados.idResumoReagente).get()).data()!;
  const estadoFisico: "SOLIDO" | "LIQUIDO" = resumo.estado_fisico;
  const densidade: number | null = especData.densidade ?? null;

  // C5: bifurcação explícita por estado físico
  let pesoVazio: number;
  let conteudoNominal: number; // g para sólidos, mL para líquidos

  if (estadoFisico === "SOLIDO") {
    // Sólidos: tudo em gramas; densidade não é usada para cálculo de volume
    if (dados.modalidade === "CONHECE_TARA") {
      if (dados.pesoFrascoVazioInformado == null)
        throw new HttpsError("invalid-argument", "CONHECE_TARA exige pesoFrascoVazioInformado.");
      pesoVazio = dados.pesoFrascoVazioInformado;
      conteudoNominal = dados.pesoTotalBalanca - pesoVazio;
    } else { // ESTIMA_MASSA
      if (dados.massaAtualEstimada == null)
        throw new HttpsError("invalid-argument", "ESTIMA_MASSA exige massaAtualEstimada.");
      conteudoNominal = dados.massaAtualEstimada;
      pesoVazio = dados.pesoTotalBalanca - conteudoNominal;
    }
    if (pesoVazio <= 0)
      throw new HttpsError("invalid-argument", "Massa atual maior que o peso total medido na balança.");
  } else { // LIQUIDO
    if (!densidade || densidade <= 0)
      throw new HttpsError("failed-precondition", "Líquido exige densidade positiva cadastrada na especificação.");
    if (dados.modalidade === "CONHECE_TARA") {
      if (dados.pesoFrascoVazioInformado == null)
        throw new HttpsError("invalid-argument", "CONHECE_TARA exige pesoFrascoVazioInformado.");
      pesoVazio = dados.pesoFrascoVazioInformado;
      conteudoNominal = (dados.pesoTotalBalanca - pesoVazio) / densidade;
    } else { // ESTIMA_VOLUME
      if (dados.volumeAtualEstimado == null)
        throw new HttpsError("invalid-argument", "ESTIMA_VOLUME exige volumeAtualEstimado.");
      conteudoNominal = dados.volumeAtualEstimado;
      pesoVazio = dados.pesoTotalBalanca - (conteudoNominal * densidade);
    }
    if (conteudoNominal <= 0 || pesoVazio <= 0)
      throw new HttpsError("invalid-argument", "Dados de volume/peso inconsistentes com a densidade.");
  }

  const agora = new Date();
  const validade = dados.validadeAberto ? new Date(dados.validadeAberto) : null;
  const validadeEfetiva = dados.validadeDesconhecida ? null : validade;
  const jaVencido = Boolean(validadeEfetiva && validadeEfetiva <= agora);

  if (jaVencido && !dados.decisaoSeJaVencido)
    throw new HttpsError("failed-precondition", "O frasco está vencido no cadastro e exige uma decisão do gestor.");
  if (dados.decisaoSeJaVencido === "QUARENTENA" && !dados.detalheStatus)
    throw new HttpsError("invalid-argument", "A entrada em quarentena exige detalhe_status.");

  return admin.firestore().runTransaction(async (tx) => {
    const { novoCodigo } = await obterProximoCodigoFrascoTx(tx);
    const frascoRef = admin.firestore().collection("Frasco_Reagente").doc();

    tx.set(frascoRef, {
      id_almoxarifado: dados.idAlmoxarifado,
      id_lote: dados.idLote ?? null,
      id_especificacao_reagente: dados.idEspecificacaoReagente, // projeção física
      id_resumo_reagente: dados.idResumoReagente,
      eh_higroscopico: resumo.eh_higroscopico, // snapshot servidor imutável
      codigo_frasco: novoCodigo,
      conteudo_nominal: conteudoNominal,   // g (sólido) ou mL (líquido)
      peso_no_cadastrado: dados.pesoTotalBalanca,
      peso_atual: dados.pesoTotalBalanca,
      peso_frasco_vazio: pesoVazio,
      medida_usada: 0,
      data_abertura: dados.aberturaHistoricaDesconhecida ? null : dados.dataAbertura,
      abertura_historica_desconhecida: dados.aberturaHistoricaDesconhecida,
      estado_fisico_frasco: "ABERTO",
      disponibilidade: "DISPONIVEL",
      validade_efetiva: validadeEfetiva,
      validade_desconhecida: dados.validadeDesconhecida ?? false,
      vencido: jaVencido,
      em_quarentena: jaVencido && dados.decisaoSeJaVencido === "QUARENTENA",
      uso_vencido_autorizado: jaVencido && dados.decisaoSeJaVencido === "DISPONIVEL",
      detalhe_status: jaVencido
        ? (dados.detalheStatus ?? (dados.decisaoSeJaVencido === "PENDENTE_DE_DESCARTE"
            ? "Frasco vencido aguardando processo institucional de descarte." : null))
        : null,
      cadastrado_em: admin.firestore.FieldValue.serverTimestamp(),
      cadastrado_por: request.auth!.uid,
    });

    return { idFrasco: frascoRef.id, codigoFrasco: novoCodigo, conteudoNominal, pesoVazio, venceuNoCadastro: jaVencido };
  });
});
\end{lstlisting}

\begin{lstlisting}[language=TypeScript, title={Abertura e cálculo de validade efetiva no servidor}]
function calcularValidadeEfetivaNaAbertura(frasco: any, dataAbertura: Date): Date | null {
  if (frasco.validade_desconhecida) return null;

  const validadeFechado = frasco.validade_fechado
    ? (typeof frasco.validade_fechado.toDate === "function"
      ? frasco.validade_fechado.toDate()
      : new Date(frasco.validade_fechado))
    : null;
  const diasDepoisDeAberto = frasco.validade_apos_aberto_dias;
  const validadeDepoisDaAbertura = diasDepoisDeAberto != null
    ? new Date(dataAbertura.getTime() + diasDepoisDeAberto * 24 * 60 * 60 * 1000)
    : null;

  if (validadeFechado && validadeDepoisDaAbertura) {
    return validadeFechado <= validadeDepoisDaAbertura ? validadeFechado : validadeDepoisDaAbertura;
  }
  return validadeFechado ?? validadeDepoisDaAbertura;
}

interface AberturaFrasco {
  idFrasco: string;
  destinoSeVencerNaAbertura?: "QUARENTENA" | "PENDENTE_DE_DESCARTE" | "DISPONIVEL";
  detalheStatus?: string;
}

export const registrarAberturaFrasco = onCall(async (request) => {
  const dados = request.data as AberturaFrasco;
  await validarPermissao(request, ["Chefe_Geral", "Gestor_Almoxarifado"]);

  const frascoRef = admin.firestore().collection("Frasco_Reagente").doc(dados.idFrasco);
  return admin.firestore().runTransaction(async (tx) => {
    const snap = await tx.get(frascoRef);
    if (!snap.exists) throw new HttpsError("not-found", "Frasco não encontrado.");
    const frasco = snap.data()!;
    await validarGestorDoAlmoxarifado(request.auth!.uid, request.auth!.token, frasco.id_almoxarifado);

    if (frasco.estado_fisico_frasco !== "FECHADO") {
      throw new HttpsError("failed-precondition", "Somente um frasco FECHADO pode ser aberto por esta operação.");
    }

    const agora = new Date();
    const validadeEfetiva = calcularValidadeEfetivaNaAbertura(frasco, agora);
    const venceu = Boolean(validadeEfetiva && validadeEfetiva <= agora);

    if (venceu && !dados.destinoSeVencerNaAbertura) {
      throw new HttpsError("failed-precondition", "A abertura tornou o frasco vencido. O gestor deve escolher o destino do frasco.");
    }
    if (venceu && dados.destinoSeVencerNaAbertura === "QUARENTENA" && !dados.detalheStatus) {
      throw new HttpsError("invalid-argument", "A entrada em quarentena exige detalhe_status.");
    }

    tx.update(frascoRef, {
      estado_fisico_frasco: "ABERTO",
      data_abertura: agora,
      validade_efetiva: validadeEfetiva,
      vencido: venceu,
      em_quarentena: venceu && dados.destinoSeVencerNaAbertura === "QUARENTENA",
      uso_vencido_autorizado: venceu && dados.destinoSeVencerNaAbertura === "DISPONIVEL",
      detalhe_status: venceu
        ? (dados.detalheStatus ?? (dados.destinoSeVencerNaAbertura === "PENDENTE_DE_DESCARTE"
          ? "Frasco vencido aguardando processo institucional de descarte."
          : null))
        : null,
    });

    return { validadeEfetiva, venceu, destino: dados.destinoSeVencerNaAbertura ?? null };
  });
});
\end{lstlisting}

\begin{lstlisting}[language=TypeScript, title={Registrar Retirada (Empréstimo)}]
interface RetiradaFrasco {
  idFrasco: string;
  idUsuarioRetirou: string;
  idLocalUsado: string;
  pesoSaida: number;
  dataDevolucaoPrevista: string;
  confirmarUsoVencido?: boolean;
  abrirNoEmprestimo?: boolean;
  finalidadeUso: "AULA_PRATICA" | "DEMONSTRACAO" | "PESQUISA_TCC_POS" | "ESTUDO_DEGRADACAO_RESIDUOS";
  justificativaMetodologica?: string;
  idAceiteTcr?: string; // evento prévio autenticado pelo retirante
  justificativaAutoAtendimento?: string;
}

export const registrarRetirada = onCall(async (request) => {
  const dados = request.data as RetiradaFrasco;
  await validarPermissao(request, ["Chefe_Geral", "Gestor_Almoxarifado"], true);

  // Verificação do papel do tomador: este UID é de outro usuário
  // (não do requisitante), portanto Custom Claims não estão disponíveis.
  // Leitura Firestore é necessária aqui (1 leitura, não 6).
  const colecoesBorrower = ["Professor", "Bolsista"];
  const [profSnap, bolsSnap] = await Promise.all([
    admin.firestore().collection("Professor").doc(dados.idUsuarioRetirou).get(),
    admin.firestore().collection("Bolsista").doc(dados.idUsuarioRetirou).get(),
  ]);
  if (!profSnap.exists && !bolsSnap.exists) {
    throw new HttpsError("failed-precondition", "Usuário selecionado não é Professor nem Bolsista.");
  }

  const frascoRef = admin.firestore().collection("Frasco_Reagente").doc(dados.idFrasco);

  return admin.firestore().runTransaction(async (tx) => {
    const frascoSnap = await tx.get(frascoRef);
    if (!frascoSnap.exists) throw new HttpsError("not-found", "Frasco não encontrado.");
    const frasco = frascoSnap.data()!;

    await validarGestorDoAlmoxarifado(request.auth!.uid, request.auth!.token, frasco.id_almoxarifado);
    if (frasco.disponibilidade !== "DISPONIVEL") {
      throw new HttpsError("failed-precondition", "Frasco não está disponível para retirada.");
    }
    if (frasco.em_quarentena) {
      throw new HttpsError("failed-precondition", "Frasco em quarentena não pode ser retirado.");
    }

    const agora = new Date();
    let frascoVencido = frasco.vencido;
    let validadeEfetiva = frasco.validade_efetiva ? frasco.validade_efetiva.toDate() : null;

    if (dados.abrirNoEmprestimo && frasco.estado_fisico_frasco === "FECHADO") {
      validadeEfetiva = calcularValidadeEfetivaNaAbertura(frasco, agora);
      frascoVencido = Boolean(validadeEfetiva && validadeEfetiva <= agora);

    }

    const excepcional = frascoVencido || frasco.validade_desconhecida;
    if (excepcional && (!frasco.uso_vencido_autorizado || !dados.confirmarUsoVencido))
      throw new HttpsError("failed-precondition", "Uso excepcional exige autorização e ciência.");
    const finalidade = dados.finalidadeUso; // nunca substituir finalidade declarada
    const pesquisa = ["PESQUISA_TCC_POS", "ESTUDO_DEGRADACAO_RESIDUOS"].includes(finalidade);
    // Helpers contratuais: todas as leituras devem ocorrer antes da primeira escrita.
    // Aceite foi criado por endpoint autenticado do retirante, não pelo gestor.
    const aceite = excepcional && pesquisa
      ? await validarAceiteTcrTx(tx, dados.idAceiteTcr, {
          retirante: dados.idUsuarioRetirou, frasco: dados.idFrasco,
          finalidade, justificativa: dados.justificativaMetodologica,
        }) : null;
    const autoAtendimento = await validarAutoAtendimentoTx(tx, request.auth!.uid,
      dados.idUsuarioRetirou, frasco.id_almoxarifado, dados.justificativaAutoAtendimento);
    const quimica = await resolverSnapshotQuimicoTx(tx, frasco);

    const emprestimoRef = admin.firestore().collection("Emprestimo_Reagente").doc();
    tx.set(emprestimoRef, {
      id_frasco_reagente: dados.idFrasco,
      id_usuario_retirou: dados.idUsuarioRetirou,
      id_almoxarifado: frasco.id_almoxarifado,
      status: "EM_USO",
      data_retirada: admin.firestore.FieldValue.serverTimestamp(),
      data_devolucao_prevista: new Date(dados.dataDevolucaoPrevista),
      id_local_usado: dados.idLocalUsado,
      peso_saida: dados.pesoSaida,
      uso_vencido_aceito: Boolean(frascoVencido),
      finalidade_uso: finalidade,
      auto_atendimento: autoAtendimento,
      justificativa_metodologica: dados.justificativaMetodologica ?? null,
      tcr_versao: aceite?.versao ?? null,
      tcr_aceito_em: aceite?.aceito_em ?? null,
      tcr_aceito_por: aceite?.uid ?? null,
      tcr_auditoria_id: aceite?.auditoria_id ?? null,
      unidade_medida_utilizada: quimica.unidade,
      densidade_aplicada: quimica.densidade,

    });

    tx.update(frascoRef, {
      disponibilidade: "EMPRESTADO",
      ...(dados.abrirNoEmprestimo && frasco.estado_fisico_frasco === "FECHADO"
        ? { estado_fisico_frasco: "ABERTO", data_abertura: agora,
            validade_efetiva: validadeEfetiva, vencido: frascoVencido } : {}),
    });
    if (aceite) vincularAceiteEAuditoriaTx(tx, aceite, emprestimoRef.id);
    if (autoAtendimento) registrarNotificacaoChefiaTx(tx, emprestimoRef.id);

    const historicoRef = admin.firestore().collection("Historico_Frasco_Reagente").doc();
    tx.set(historicoRef, {
      id_frasco_reagente: dados.idFrasco,
      id_almoxarifado: frasco.id_almoxarifado,
      id_gestor: request.auth!.uid,
      tipo: "SAIU",
      id_emprestimo_reagente: emprestimoRef.id,
      timestamp: admin.firestore.FieldValue.serverTimestamp(),
    });

    return { idEmprestimo: emprestimoRef.id };
  });
});
\end{lstlisting}

\begin{lstlisting}[language=TypeScript, title={Registrar Devolução com atualização física e histórico completo}]
interface DevolucaoFrasco {
  idEmprestimo: string;
  pesoRetorno: number;
  destinoPosDevolucao?: "QUARENTENA" | "PENDENTE_DE_DESCARTE" | "DISPONIVEL";
}

export const registrarDevolucao = onCall(async (request) => {
  const dados = request.data as DevolucaoFrasco;
  await validarPermissao(request, ["Chefe_Geral", "Gestor_Almoxarifado"], true);

  const emprestimoRef = admin.firestore().collection("Emprestimo_Reagente").doc(dados.idEmprestimo);

  return admin.firestore().runTransaction(async (tx) => {
    const emprestimoSnap = await tx.get(emprestimoRef);
    if (!emprestimoSnap.exists) throw new HttpsError("not-found", "Empréstimo não encontrado.");
    const emprestimo = emprestimoSnap.data()!;
    if (emprestimo.status === "DEVOLVIDO" || emprestimo.status === "DEVOLVIDO_COM_ATRASO") {
      throw new HttpsError("failed-precondition", "Empréstimo já foi devolvido.");
    }

    const frascoRef = admin.firestore().collection("Frasco_Reagente").doc(emprestimo.id_frasco_reagente);
    const frascoSnap = await tx.get(frascoRef);
    const frasco = frascoSnap.data()!;
    await validarGestorDoAlmoxarifado(request.auth!.uid, request.auth!.token, frasco.id_almoxarifado);

    const liquido = emprestimo.unidade_medida_utilizada === "ml";
    const densidade = liquido ? emprestimo.densidade_aplicada : null;
    if (liquido && (!Number.isFinite(densidade) || densidade <= 0))
      throw new HttpsError("failed-precondition", "Densidade histórica inválida.");
    const pesoConsumido = Math.max(0, emprestimo.peso_saida - dados.pesoRetorno);

    if (typeof frasco.eh_higroscopico !== "boolean")
      throw new HttpsError("failed-precondition", "Snapshot histórico ausente: reconciliar frasco.");
    const deltaMax = frasco.eh_higroscopico
      ? Math.max(2.0, 0.02 * emprestimo.peso_saida)
      : Math.max(1.0, 0.005 * emprestimo.peso_saida);
    const limiteMaxRetorno = emprestimo.peso_saida + deltaMax;
    if (!Number.isFinite(dados.pesoRetorno) || dados.pesoRetorno <= 0 || dados.pesoRetorno > limiteMaxRetorno)
      throw new HttpsError("invalid-argument", "Peso fora da tolerância híbrida Q06.");
    const avisoHigroscopico = dados.pesoRetorno > emprestimo.peso_saida;
    const volumeUtilizado = densidade ? pesoConsumido / densidade : pesoConsumido;

    const agora = new Date();
    // Normalização: Prazo encerra-se às 23:59:59.999 do dia previsto
    const prazoLimite = new Date(emprestimo.data_devolucao_prevista.toDate());
    prazoLimite.setHours(23, 59, 59, 999);
    const atrasado = agora > prazoLimite;

    const validadeEfetiva = frasco.validade_efetiva?.toDate?.() ?? frasco.validade_efetiva ?? null;
    const venceuAgora = Boolean(validadeEfetiva && validadeEfetiva <= agora);
    const frascoVencido = Boolean(frasco.vencido || venceuAgora);

    const atualizacaoFrasco: Record<string, unknown> = {
      disponibilidade: "DISPONIVEL",
      data_ultima_pesagem: admin.firestore.FieldValue.serverTimestamp(),
      peso_atual: dados.pesoRetorno,
      medida_usada: admin.firestore.FieldValue.increment(pesoConsumido), // sempre g
      vencido: frascoVencido,
    };

    if (frascoVencido && dados.destinoPosDevolucao === "QUARENTENA") {
      atualizacaoFrasco.em_quarentena = true;
      atualizacaoFrasco.uso_vencido_autorizado = false;
      atualizacaoFrasco.detalhe_status = "Frasco vencido retido em quarentena após devolução.";
    } else if (frascoVencido && dados.destinoPosDevolucao === "DISPONIVEL") {
      atualizacaoFrasco.em_quarentena = false;
      atualizacaoFrasco.uso_vencido_autorizado = true;
    }

    tx.update(emprestimoRef, {
      status: atrasado ? "DEVOLVIDO_COM_ATRASO" : "DEVOLVIDO",
      data_devolucao_efetuada: admin.firestore.FieldValue.serverTimestamp(),
      id_usuario_devolveu: request.auth!.uid,
      peso_retorno: dados.pesoRetorno,
      medida_utilizada: volumeUtilizado,
      unidade_medida_utilizada: densidade ? "ml" : "g",
    ...(avisoHigroscopico && { aviso: "higroscopico_suspeito" }),
    });
    tx.update(frascoRef, atualizacaoFrasco);

    if (avisoHigroscopico) {
      tx.set(admin.firestore().collection("Historico_Frasco_Reagente").doc(), {
        id_frasco_reagente: frascoRef.id,
        id_almoxarifado: frasco.id_almoxarifado,
        id_gestor: request.auth!.uid,
        id_emprestimo_reagente: emprestimoRef.id,
        tipo: "AJUSTE", campo_ajustado: "ganho_massa_higroscopia",
        peso_anterior: emprestimo.peso_saida, peso_novo: dados.pesoRetorno,
        medida_ajustada: dados.pesoRetorno - emprestimo.peso_saida,
        unidade_medida_ajustada: "g",
        timestamp: admin.firestore.FieldValue.serverTimestamp(),
      });
    }

    // Registro obrigatório de devolução no histórico
    const histRef = admin.firestore().collection("Historico_Frasco_Reagente").doc();
    tx.set(histRef, {
      id_frasco_reagente: emprestimo.id_frasco_reagente,
      id_almoxarifado: frasco.id_almoxarifado,
      id_gestor: request.auth!.uid,
      tipo: "ENTROU",
      id_emprestimo_reagente: emprestimoRef.id,
      peso_anterior: emprestimo.peso_saida,
      peso_novo: dados.pesoRetorno,
      medida_ajustada: volumeUtilizado,
      unidade_medida_ajustada: densidade ? "ml" : "g",
      timestamp: admin.firestore.FieldValue.serverTimestamp(),
    });

    return { pesoConsumido, volumeUtilizado, atrasado, frascoVencido,
           avisoHigroscopico: avisoHigroscopico || false };
  });
});
\end{lstlisting}

\subsubsection{Validação síncrona de validade no cadastro e na abertura}
A validade não depende apenas do job diário. No cadastro do frasco, a Cloud Function calcula \textit{validade\_efetiva} e compara esse valor com o relógio do servidor. Se a validade já expirou, a criação só pode prosseguir com uma decisão explícita do gestor: \texttt{QUARENTENA}, \texttt{PENDENTE\_DE\_DESCARTE} ou \texttt{DISPONIVEL} com \textit{uso\_vencido\_autorizado}. A primeira opção exige \textit{detalhe\_status}. A terceira representa uma exceção operacional deliberada: reagentes vencidos podem continuar sendo utilizados em práticas menos exigentes ou demonstrações quando o gestor registra explicitamente essa autorização.

Na abertura do frasco, a função \textit{calcularValidadeEfetivaNaAbertura} aplica:
\[
\textit{validade\_efetiva} = \min(\textit{validade\_fechado},\; \textit{data\_abertura} + \textit{validade\_apos\_aberto\_dias}),
\]
quando os dois limites existem.

\subsubsection{Jobs Agendados: Vencimento, Atraso e Materialização}

\begin{lstlisting}[language=TypeScript, title={Vencimento e Atraso com busca preventiva de devoluções}]
import { onSchedule } from "firebase-functions/v2/scheduler";
import { onDocumentCreated, onDocumentDeleted, onDocumentUpdated } from "firebase-functions/v2/firestore";

export const verificarVencimentosEAtrasos = onSchedule("every day 03:00", async () => {
  const agora = new Date();
  const inicioHoje = inicioDoDiaInstitucional(agora);
  const fimDeAmanha = fimDoDiaInstitucional(adicionarDias(inicioHoje, 1));

  const emprestimosNaJanela = await admin.firestore().collection("Emprestimo_Reagente")
    .where("status", "==", "EM_USO")
    .where("data_devolucao_prevista", "<=", fimDeAmanha)
    .get();

  const batchAtrasados = admin.firestore().batch();
  const notificacoesPreventivas: Array<{ idEmprestimo: string; idUsuario: string; quando: "HOJE" | "AMANHA" }> = [];

  emprestimosNaJanela.docs.forEach((doc) => {
    const dados = doc.data();
    const prevista = dados.data_devolucao_prevista.toDate();

    if (prevista < inicioHoje) {
      batchAtrasados.update(doc.ref, { status: "ATRASADO" });
      return;
    }

    if (mesmoDia(prevista, inicioHoje)) {
      notificacoesPreventivas.push({ idEmprestimo: doc.id, idUsuario: dados.id_usuario_retirou, quando: "HOJE" });
    } else {
      notificacoesPreventivas.push({ idEmprestimo: doc.id, idUsuario: dados.id_usuario_retirou, quando: "AMANHA" });
    }
  });

  await batchAtrasados.commit();

  for (const item of notificacoesPreventivas) {
    await notificarProfessorDevolucaoUmaVez(item.idUsuario, item.idEmprestimo, item.quando);
  }

  const totalAtrasados = emprestimosNaJanela.docs.filter((doc) => {
    const prevista = doc.data().data_devolucao_prevista.toDate();
    return prevista < inicioHoje;
  }).length;
  if (totalAtrasados > 0) {
    await notificarGestoresDeReagentes("ENTREGA_ATRASADA", totalAtrasados);
  }

  const frascosVencendo = await admin.firestore().collection("Frasco_Reagente")
    .where("vencido", "==", false)
    .where("validade_efetiva", "<=", agora)
    .get();

  const batchVencidos = admin.firestore().batch();
  frascosVencendo.docs.forEach((doc) => {
    batchVencidos.update(doc.ref, { vencido: true });
    const histRef = admin.firestore().collection("Historico_Frasco_Reagente").doc();
    batchVencidos.set(histRef, {
      id_frasco_reagente: doc.id,
      id_almoxarifado: doc.data().id_almoxarifado,
      tipo: "VENCEU",
      timestamp: admin.firestore.FieldValue.serverTimestamp(),
    });
  });
  await batchVencidos.commit();
  if (!frascosVencendo.empty) {
    await notificarGestoresDeReagentes("FRASCOS_VENCIDOS", frascosVencendo.size);
  }
});

function adicionarDias(data: Date, dias: number): Date {
  const result = new Date(data);
  result.setDate(result.getDate() + dias);
  return result;
}

function mesmoDia(a: Date, b: Date): boolean {
  return a.getFullYear() === b.getFullYear() &&
         a.getMonth() === b.getMonth() &&
         a.getDate() === b.getDate();
}

function inicioDoDiaInstitucional(data: Date): Date {
  const result = new Date(data);
  result.setHours(0, 0, 0, 0);
  return result;
}

function fimDoDiaInstitucional(data: Date): Date {
  const result = new Date(data);
  result.setHours(23, 59, 59, 999);
  return result;
}

export const onFrascoCriado = onDocumentCreated("Frasco_Reagente/{frascoId}", async (event) => {
  const idLote = event.data?.data().id_lote;
  if (idLote) await atualizarContadorLote(idLote, +1);
});

export const onFrascoRemovido = onDocumentDeleted("Frasco_Reagente/{frascoId}", async (event) => {
  const idLote = event.data?.data().id_lote;
  if (idLote) await atualizarContadorLote(idLote, -1);
});

async function atualizarContadorLote(idLote: string, delta: number) {
  await admin.firestore().collection("Lote_Materializado").doc(idLote).set(
    { id_lote: idLote, qtd_frascos_cadastrados: admin.firestore.FieldValue.increment(delta) },
    { merge: true }
  );
}
\end{lstlisting}

\subsubsection{Fluxo de Bens Patrimoniais: Requisições}

\begin{lstlisting}[language=TypeScript, title={Criação e Resposta de Requisição de Edição (C2: lock determinístico)}]
export const criarRequisicaoEdicaoBem = onCall(async (request) => {
  await validarPermissao(request, ["Professor"]);
  const dados = request.data as { idBemPatrimonial: string; novoNome?: string; novoStatus?: string; novoEstadoConservacao?: string; novoIdLocal?: string; motivo: string };

  // C2: Lock determinístico - elimina phantom reads em transações Firestore.
  // Se dois professores enviarem simultaneamente, apenas um cria o lock; o outro
  // recebe ALREADY_EXISTS da transação e falha antes de criar a requisição.
  const lockId = `bem_edicao_${dados.idBemPatrimonial}`;
  const lockRef = admin.firestore().collection("Locks_Requisicao_Patrimonio").doc(lockId);
  const reqRef  = admin.firestore().collection("Requisicao_Edicao_Bem_Patrimonial").doc();

  return admin.firestore().runTransaction(async (tx) => {
    const lockSnap = await tx.get(lockRef);
    if (lockSnap.exists) {
      throw new HttpsError("failed-precondition", "Já existe uma requisição de edição pendente para este bem.");
    }
    const bemSnap = await tx.get(admin.firestore().collection("Bem_Patrimonial").doc(dados.idBemPatrimonial));
    if (!bemSnap.exists) throw new HttpsError("not-found", "Bem não encontrado.");
    // Cria lock e requisição atomicamente
    tx.set(lockRef, { criado_em: admin.firestore.FieldValue.serverTimestamp() });
    tx.set(reqRef, {
      id_bem_patrimonial: dados.idBemPatrimonial,
      novo_nome: dados.novoNome ?? null,
      versao_bem_origem: bemSnap.data()!.versao,
      novo_id_resumo_bem_patrimonial: null,
      novo_status: dados.novoStatus ?? null,
      novo_estado_conservacao: dados.novoEstadoConservacao ?? null,
      novo_id_local: dados.novoIdLocal ?? null,
      motivo: dados.motivo,
      status: "pendente",
      feita_em: admin.firestore.FieldValue.serverTimestamp(),
      id_usuario_solicitante: request.auth!.uid,
    });
    return { idRequisicao: reqRef.id };
  });
});

export const responderRequisicaoEdicaoBem = onCall(async (request) => {
  await validarPermissao(request, ["Chefe_Geral", "Gestor_Bens_Patrimoniais"], true);
  const { idRequisicao, aprovar, justificativa, idResumoAlvo } = request.data;
  return admin.firestore().runTransaction(async (tx) => {
    const reqRef = admin.firestore().collection("Requisicao_Edicao_Bem_Patrimonial").doc(idRequisicao);
    const req = (await tx.get(reqRef)).data();
    if (!req || req.status !== "pendente") throw new HttpsError("failed-precondition", "Requisição indisponível.");
    const bemRef = admin.firestore().collection("Bem_Patrimonial").doc(req.id_bem_patrimonial);
    const bem = (await tx.get(bemRef)).data();
    if (!bem || bem.versao !== req.versao_bem_origem)
      throw new HttpsError("aborted", "Bem alterado; revisar proposta.");
    // Resolve alvo existente ou prepara criação; helper só lê nesta etapa.
    const alvo = aprovar && req.novo_nome
      ? await prepararResumoAlvoTx(tx, idResumoAlvo, req.novo_nome) : null;
    const alteracoes = calcularAlteracoesPermitidas(req, bem);
    if (alvo) {
      if (alvo.novo) tx.create(alvo.ref, alvo.dados);
      alteracoes.id_resumo_bem_patrimonial = alvo.ref.id;
      alteracoes.nome_equipamento = alvo.dados.nome;
    }
    if (aprovar) {
      tx.update(bemRef, { ...alteracoes, versao: bem.versao + 1 });
      registrarHistoricoBemTx(tx, bemRef, bem, alteracoes, request.auth!.uid);
    }
    tx.update(reqRef, {
      status: aprovar ? "aprovada" : "rejeitada",
      novo_id_resumo_bem_patrimonial: alvo?.ref.id ?? null,
      respondida_em: admin.firestore.FieldValue.serverTimestamp(),
      id_usuario_respondente: request.auth!.uid, justificativa_resposta: justificativa,
    });
    tx.delete(admin.firestore().collection("Locks_Requisicao_Patrimonio").doc(`bem_edicao_${bemRef.id}`));
    return { status: aprovar ? "aprovada" : "rejeitada" };
  });
});
\end{lstlisting}

\begin{lstlisting}[language=TypeScript, title={criarRequisicaoAdicaoBem com lock determinístico (C2)}]
export const criarRequisicaoAdicaoBem = onCall(async (request) => {
  await validarPermissao(request, ["Professor"]);

  const dados = request.data as {
    numeroPatrimonioProposto: string;
    estadoConservacaoProposto: string;
    idLocal: string;
    nomeResponsavelProposto: string;
    idResumoBemPatrimonial?: string;
    nomeResumoProposto?: string;
    descricaoResumoProposta?: string;
    motivo: string;
  };

  // C2: Lock determinístico por número de patrimônio proposto
  const lockId = `bem_adicao_${dados.numeroPatrimonioProposto}`;
  const lockRef = admin.firestore().collection("Locks_Requisicao_Patrimonio").doc(lockId);
  const reqRef  = admin.firestore().collection("Requisicao_Adicao_Bem_Patrimonial").doc();

  return admin.firestore().runTransaction(async (tx) => {
    const lockSnap = await tx.get(lockRef);
    if (lockSnap.exists) {
      throw new HttpsError("failed-precondition", "Já existe requisição pendente para este número de patrimônio.");
    }
    tx.set(lockRef, { criado_em: admin.firestore.FieldValue.serverTimestamp() });
    tx.set(reqRef, {
      numero_patrimonio_proposto: dados.numeroPatrimonioProposto,
      estado_conservacao_proposto: dados.estadoConservacaoProposto,
      photo_url_proposta: null,
      nome_responsavel_proposto: dados.nomeResponsavelProposto,
      id_local: dados.idLocal,
      id_resumo_bem_patrimonial: dados.idResumoBemPatrimonial ?? null,
      nome_resumo_proposto: dados.nomeResumoProposto ?? null,
      descricao_resumo_proposta: dados.descricaoResumoProposta ?? null,
      motivo: dados.motivo,
      status: "pendente",
      feita_em: admin.firestore.FieldValue.serverTimestamp(),
      id_usuario_solicitante: request.auth!.uid,
      respondida_em: null,
      id_usuario_respondente: null,
    });
    return { idRequisicao: reqRef.id };
  });
});

export const responderRequisicaoAdicaoBem = onCall(async (request) => {
  await validarPermissao(request, ["Chefe_Geral", "Gestor_Bens_Patrimoniais"], true);
  const { idRequisicao, aprovar, justificativa } = request.data as { idRequisicao: string; aprovar: boolean; justificativa: string };
  const reqRef = admin.firestore().collection("Requisicao_Adicao_Bem_Patrimonial").doc(idRequisicao);

  return admin.firestore().runTransaction(async (tx) => {
    const reqSnap = await tx.get(reqRef);
    if (!reqSnap.exists) throw new HttpsError("not-found", "Requisição não encontrada.");
    const req = reqSnap.data()!;
    if (req.status !== "pendente") throw new HttpsError("failed-precondition", "Requisição já respondida.");

    // C2: Remove o lock ao responder (libera possível nova requisição futura)
    const lockRef = admin.firestore().collection("Locks_Requisicao_Patrimonio")
      .doc(`bem_adicao_${req.numero_patrimonio_proposto}`);
    tx.delete(lockRef);

    let idBemCriado: string | null = null;

    if (aprovar) {
      let idResumo = req.id_resumo_bem_patrimonial;

      // Se não havia resumo existente, cria atômica e transacionalmente o Resumo_Bem_Patrimonial
      if (!idResumo && req.nome_resumo_proposto) {
        const novoResumoRef = admin.firestore().collection("Resumo_Bem_Patrimonial").doc();
        tx.set(novoResumoRef, {
          nome: req.nome_resumo_proposto,
          descricao: req.descricao_resumo_proposta ?? "",
          criado_em: admin.firestore.FieldValue.serverTimestamp(),
        });
        idResumo = novoResumoRef.id;
      }

      if (!idResumo) throw new HttpsError("failed-precondition", "A aprovação exige vincular um resumo.");

      const bemRef = admin.firestore().collection("Bem_Patrimonial").doc();
      idBemCriado = bemRef.id;

      tx.set(bemRef, {
        id_resumo_bem_patrimonial: idResumo,
        nome_equipamento: req.nome_resumo_proposto ?? "Equipamento",
        numero_patrimonio: req.numero_patrimonio_proposto,
        estado_conservacao: req.estado_conservacao_proposto,
        id_local: req.id_local,
        photo_url: req.photo_url_proposta ?? "https://placeholder",
        documento_dado_baixa_pdf_url: null,
        nome_responsavel_sei: req.nome_responsavel_proposto,
        status: "Ativo",
        cadastrado_em: admin.firestore.FieldValue.serverTimestamp(),
      });

      tx.update(reqRef, {
        status: "aprovada",
        respondida_em: admin.firestore.FieldValue.serverTimestamp(),
        id_usuario_respondente: request.auth!.uid,
        id_bem_patrimonial_se_aprovado: idBemCriado,
        justificativa_resposta: justificativa,
      });
    } else {
      tx.update(reqRef, {
        status: "rejeitada",
        respondida_em: admin.firestore.FieldValue.serverTimestamp(),
        id_usuario_respondente: request.auth!.uid,
        justificativa_resposta: justificativa,
      });
    }

    return { status: aprovar ? "aprovada" : "rejeitada", idBemCriado };
  });
});

export const onResumoBemPatrimonialNomeAtualizado = onDocumentUpdated(
  "Resumo_Bem_Patrimonial/{resumoId}",
  async (event) => {
    const antes = event.data?.before.data();
    const depois = event.data?.after.data();
    if (!antes || !depois || antes.nome === depois.nome) return;

    const bens = await admin.firestore().collection("Bem_Patrimonial")
      .where("id_resumo_bem_patrimonial", "==", event.params.resumoId)
      .get();

    const batch = admin.firestore().batch();
    bens.docs.forEach((bem) => batch.update(bem.ref, {
      nome_equipamento: depois.nome,
    }));
    await batch.commit();
  }
);

export const onLocalAtualizado = onDocumentUpdated(
  "Local/{localId}",
  async (event) => {
    const antes = event.data?.before.data();
    const depois = event.data?.after.data();
    if (!antes || !depois) return;
    if (antes.predio === depois.predio && antes.andar === depois.andar && antes.sala === depois.sala) return;

    const bens = await admin.firestore().collection("Bem_Patrimonial")
      .where("id_local", "==", event.params.localId)
      .get();

    const batch = admin.firestore().batch();
    bens.docs.forEach((bem) => batch.update(bem.ref, {
      predio: depois.predio,
      andar: depois.andar,
      sala: depois.sala,
    }));
    await batch.commit();
  }
);

export const ingressarEmTurmaPorCodigo = onCall(async (request) => {
  await validarPermissao(request, ["Aluno"]);
  const { codigoTurma } = request.data as { codigoTurma: string };

  const turmaSnap = await admin.firestore().collection("Turma")
    .where("codigo_turma", "==", codigoTurma)
    .where("status", "==", "Ativo")
    .limit(1).get();
  if (turmaSnap.empty) throw new HttpsError("not-found", "Turma não encontrada ou arquivada.");
  const turmaDoc = turmaSnap.docs[0];
  const turma = turmaDoc.data();

  return admin.firestore().runTransaction(async (tx) => {
    // Checar capacidade
    const alunosSnap = await tx.get(
      turmaDoc.ref.collection("Alunos")
    );
    if (alunosSnap.size >= turma.capacidade) {
      throw new HttpsError("failed-precondition", "Turma lotada.");
    }

    // Checar se já está matriculado
    const jaMatriculado = alunosSnap.docs.some(
      (d) => d.id === request.auth!.uid
    );
    if (jaMatriculado) {
      throw new HttpsError("already-exists", "Você já está nesta turma.");
    }

    // P14: Checar se foi removido anteriormente
    const historico = await tx.get(
      turmaDoc.ref.collection("HistoricoAlunos")
        .where("id_aluno", "==", request.auth!.uid)
        .where("tipo", "==", "exclusao_aluno")
        .limit(1)
    );
    if (!historico.empty) {
      throw new HttpsError("failed-precondition",
        "Você foi removido desta turma pelo professor. O re-ingresso requer convite explícito.");
    }

    // Matricular
    tx.set(turmaDoc.ref.collection("Alunos").doc(request.auth!.uid), {
      id_aluno: request.auth!.uid,
      ingressou_em: admin.firestore.FieldValue.serverTimestamp(),
    });
    tx.set(turmaDoc.ref.collection("HistoricoAlunos").doc(), {
      id_aluno: request.auth!.uid,
      tipo: "inclusao_aluno",
      timestamp: admin.firestore.FieldValue.serverTimestamp(),
    });

    return { idTurma: turmaDoc.id, nomeTurma: turma.nome_turma };
  });
});
\end{lstlisting}

\subsubsection{Relatórios em PDF}
Pseudo-códigos das funções de relatório, revisados: os nomes de campo passam a corresponder exatamente ao modelo (\textit{data\_devolucao\_efetuada}, não \textit{data\_devolucao}; filtro por \textit{id\_almoxarifado} denormalizado em \textit{Emprestimo\_Reagente}, não por um inexistente \textit{id\_local\_saida}), e \textit{gerarRelatorioAlmoxarifado} passa a checar o escopo do gestor, não apenas o papel.

\begin{lstlisting}[language=TypeScript, title={Configuração Base e Relatório de Almoxarifado}]
import PDFDocument from "pdfkit-table";

interface FiltrosAlmoxarifado {
  idAlmoxarifado: string;
  mes: number;
  ano: number;
}

export const gerarRelatorioAlmoxarifado = onCall(async (request) => {
  const { idAlmoxarifado, mes, ano } = request.data as FiltrosAlmoxarifado;
  await validarPermissao(request, ["Chefe_Geral", "Gestor_Almoxarifado"]);
  await validarGestorDoAlmoxarifado(request.auth!.uid, request.auth!.token, idAlmoxarifado);

  const dataInicio = new Date(ano, mes - 1, 1);
  const dataFim = new Date(ano, mes, 0);

  const devolucoesSnapshot = await admin.firestore().collection("Emprestimo_Reagente")
    .where("id_almoxarifado", "==", idAlmoxarifado)
    .where("data_devolucao_efetuada", ">=", dataInicio)
    .where("data_devolucao_efetuada", "<=", dataFim)
    .get();

  const volumeTotalConsumido = devolucoesSnapshot.docs.reduce(
    (acc, doc) => acc + (doc.data().medida_utilizada || 0), 0
  );

  const historicoSnapshot = await admin.firestore().collection("Historico_Frasco_Reagente")
    .where("id_almoxarifado", "==", idAlmoxarifado)
    .where("timestamp", ">=", dataInicio)
    .where("timestamp", "<=", dataFim)
    .orderBy("timestamp", "asc").get();

  const doc = new PDFDocument({ margin: 30, size: "A4" });
  const buffer = await buildPdfBuffer(doc, historicoSnapshot.docs, volumeTotalConsumido);

  const fileName = `relatorios/almoxarifado_${idAlmoxarifado}_${mes}_${ano}_${Date.now()}.pdf`;
  const fileRef = admin.storage().bucket().file(fileName);
  await fileRef.save(buffer, { contentType: "application/pdf" });

  const [url] = await fileRef.getSignedUrl({ action: "read", expires: Date.now() + 3600000 });
  return { url };
});
\end{lstlisting}

\begin{lstlisting}[language=TypeScript, title={Relatório Mensal de um Prédio (Filtros Compostos)}]
interface FiltrosPredio {
  predio: string;
  mes: number;
  ano: number;
  andar?: string;
  sala?: string;
  estadoConservacao?: string;
  status?: string;
}

export const gerarRelatorioBensPredio = onCall(async (request) => {
  const filtros = request.data as FiltrosPredio;
  await validarPermissao(request, ["Chefe_Geral", "Gestor_Bens_Patrimoniais"]);

  let query = admin.firestore().collection("Bem_Patrimonial").where("predio", "==", filtros.predio);
  if (filtros.andar) query = query.where("andar", "==", filtros.andar);
  if (filtros.sala) query = query.where("sala", "==", filtros.sala);
  if (filtros.estadoConservacao) query = query.where("estado_conservacao", "==", filtros.estadoConservacao);
  if (filtros.status) query = query.where("status", "==", filtros.status);

  const bensSnapshot = await query.get();

  const doc = new PDFDocument({ size: "A4" });
  const buffer = await buildPredioPdfBuffer(doc, bensSnapshot.docs, filtros);

  const fileRef = admin.storage().bucket().file(`relatorios/predio_${filtros.predio}_${Date.now()}.pdf`);
  await fileRef.save(buffer, { contentType: "application/pdf" });

  const [url] = await fileRef.getSignedUrl({ action: "read", expires: Date.now() + 3600000 });
  return { url };
});
\end{lstlisting}

\begin{lstlisting}[language=TypeScript, title={Relatório Geral e Período Personalizado}]
interface FiltrosGeralEPersonalizado {
  dataInicio: string;
  dataFim: string;
  entidade: "Bens_Patrimoniais" | "Reagentes";
  prediosSelecionados?: string[];
}

export const gerarRelatorioPersonalizado = onCall(async (request) => {
  const { dataInicio, dataFim, entidade, prediosSelecionados } = request.data as FiltrosGeralEPersonalizado;
  await validarPermissao(request, ["Chefe_Geral", "Gestor_Bens_Patrimoniais", "Gestor_Almoxarifado"]);

  const dataIniObj = new Date(dataInicio);
  const dataFimObj = new Date(dataFim);
  let snapshot;

  if (entidade === "Bens_Patrimoniais") {
    let query = admin.firestore().collection("Historico_Bem_Patrimonial")
      .where("timestamp", ">=", dataIniObj)
      .where("timestamp", "<=", dataFimObj);
    if (prediosSelecionados && prediosSelecionados.length > 0) {
      query = query.where("predio", "in", prediosSelecionados);
    }
    snapshot = await query.get();
  } else {
    snapshot = await admin.firestore().collection("Emprestimo_Reagente")
      .where("data_devolucao_efetuada", ">=", dataIniObj)
      .where("data_devolucao_efetuada", "<=", dataFimObj)
      .get();
  }

  const doc = new PDFDocument({ layout: "landscape", size: "A4" });
  const buffer = await buildPersonalizadoPdfBuffer(doc, snapshot.docs, entidade);

  const fileRef = admin.storage().bucket().file(`relatorios/personalizado_${Date.now()}.pdf`);
  await fileRef.save(buffer, { contentType: "application/pdf" });
  const [url] = await fileRef.getSignedUrl({ action: "read", expires: Date.now() + 3600000 });
  return { url };
});
\end{lstlisting}

\begin{lstlisting}[language=TypeScript, title={Módulo de Etiquetas em PDF: Virgens e Reimpressão}]
// ============================================================================
// GERACAO DE ETIQUETAS VIRGENS (NOVOS FRASCOS) - ABA 1
// ============================================================================
interface DadosEtiquetasVirgens {
  codigoInicial: number;
  codigoFinal: number;
  startRow?: number; // 1 a 10 (Grid A4)
  startCol?: number; // 1 a 3
}

export const gerarPdfEtiquetasVirgens = onCall(async (request) => {
  await validarPermissao(request, ["Chefe_Geral", "Gestor_Almoxarifado"]);
  const dados = request.data as DadosEtiquetasVirgens;

  const total = dados.codigoFinal - dados.codigoInicial + 1;
  if (total <= 0 || total > 50) {
    throw new HttpsError("invalid-argument", "O lote deve conter entre 1 e 50 etiquetas por impressão.");
  }

  // Audita a sessao de geracao de etiquetas virgens (Secao 4.13)
  await admin.firestore().collection("Impressao_Etiqueta_Frasco").add({
    gerado_em: admin.firestore.FieldValue.serverTimestamp(),
    gerado_por: request.auth!.uid,
    codigo_inicial: dados.codigoInicial,
    codigo_final: dados.codigoFinal,
  });

  const codigos = Array.from({ length: total }, (_, i) => `LCQUI-${dados.codigoInicial + i}`);
  const buffer = await gerarPdfEtiquetasNovosFrascos(codigos, dados.startRow || 1, dados.startCol || 1);

  const fileName = `etiquetas/virgens_${dados.codigoInicial}_a_${dados.codigoFinal}_${Date.now()}.pdf`;
  const fileRef = admin.storage().bucket().file(fileName);
  await fileRef.save(buffer, { contentType: "application/pdf" });
  const [url] = await fileRef.getSignedUrl({ action: "read", expires: Date.now() + 3600000 });

  return { url };
});

// ============================================================================
// REIMPRESSAO E FICHA DE CONFERENCIA (FRASCOS JA CADASTRADOS) - ABA 2
// ============================================================================
interface DadosReimpressao {
  frascoIds: string[]; // Máximo 10 frascos
}

export const gerarPdfReimpressaoFrascos = onCall(async (request) => {
  await validarPermissao(request, ["Chefe_Geral", "Gestor_Almoxarifado"]);
  const { frascoIds } = request.data as DadosReimpressao;

  // Regra 7.2.10: Teto rigido de seguranca quimica
  if (!frascoIds || frascoIds.length === 0 || frascoIds.length > 10) {
    throw new HttpsError("invalid-argument", "Selecione entre 1 e 10 frascos por sessao de reimpressao.");
  }

  // 1. Busca dados completos, historico recente e ultimo emprestimo dos frascos
  const dadosFrascos = await buscarDadosCompletosParaConferencia(frascoIds);

  // 2. Registra auditoria individual de 2a via para cada frasco (Secao 4.41)
  const batch = admin.firestore().batch();
  for (const frascoId of frascoIds) {
    const auditRef = admin.firestore().collection("Registro_de_Auditoria").doc();
    batch.set(auditRef, {
      id_usuario: request.auth!.uid,
      acao: "REIMPRESSAO_ETIQUETA",
      tipo_entidade_sofre_acao: "FRASCO_REAGENTE",
      id_do_objeto_da_entidade: frascoId,
      acao_feita_em: admin.firestore.FieldValue.serverTimestamp(),
      metadata: { motivo: "segunda_via_conferencia" },
    });
  }
  await batch.commit();

  // 3. Gera PDF com Ficha de Conferencia + 1 Etiqueta Centralizada por frasco
  const buffer = await gerarPdfFrascosCadastrados(dadosFrascos);

  const fileName = `etiquetas/reimpressao_${Date.now()}.pdf`;
  const fileRef = admin.storage().bucket().file(fileName);
  await fileRef.save(buffer, { contentType: "application/pdf" });
  const [url] = await fileRef.getSignedUrl({ action: "read", expires: Date.now() + 3600000 });

  return { url };
});
\end{lstlisting}



\subsection{Consolidação do planejamento complementar e contrato de execução}
\label{sec:consolidacao}
O inventário complementar inclui Planejamento\_Etiquetas\_LCQUI.md, REGRAS\_MULTI\_ROLE.md, COMPILACAO\_NIX\_LCQUI.md, README.md e os quatro documentos de acompanhamento. Os README de scaffold do frontend não adicionam requisitos de negócio; instruções AGENTS/CLAUDE orientam desenvolvimento, não a especificação do produto. O novo arquivo DUVIDAS\_DOCUMENTACAO\_LCQUI.md registra Q01--Q14 com três alternativas cada; a opção recomendada não é aprovação.

O planejamento de etiquetas foi incorporado às seções 4 (auditoria), 5.9 (contador e campos), 7 (regras), 8 (UI-08), 9 (ALM-07), 10 (geração) e 11 (acesso). O documento de Multi-Role foi incorporado integralmente pelos IDs RN-ROLE-01 a RN-ROLE-15 na seção 7. Os Markdown de origem permanecem como histórico; esta versão explicita as decisões conflitantes em vez de criar fontes concorrentes.

\paragraph{Etiquetas e geometria.}
Papel A4 210 por 297 mm; três colunas e dez linhas; etiqueta 65 por 26,5 mm; margens laterais 7,5 mm e superior/inferior 16 mm; espaçamentos horizontal/vertical zero. Linhas cinza contínuas orientam corte. Offset é aplicado apenas na primeira página; ao completar a grade, reiniciar na célula 1/1 da seguinte. Etiqueta virgem contém código, espaço manuscrito para data/reagente e Code 128. Segunda via contém nome e data cadastrados e uma etiqueta centralizada por ficha. Gerar barras em memória com bwip-js e PDF com pdfkit; preservar proporção e área livre ao redor do código, testar leitura física e códigos longos. Casos mínimos: uma etiqueta, coluna 3, offset 4/2, última célula 10/3, 30/36 itens, nome longo, frasco sem histórico e rejeição de 11 IDs na segunda via. Os testes relatados no planejamento são evidência histórica de protótipo, não homologação da aplicação atual.

\paragraph{Contador e auditoria.}
Adota-se o caminho já usado pelo sistema, Contador\_Codigo\_Frasco/singleton e ultimo\_codigo\_gerado. Contadores/codigo\_frasco no planejamento é proposta antiga, não outro contador ativo. Cadastro incrementa contador e cria frasco atomicamente; impressão não reserva código. Lotes virgens usam Impressao\_Etiqueta\_Frasco; segunda via usa Registro\_de\_Auditoria. A sugestão contraditória de armazenar segunda via também em Impressao não é adotada. Custos reais incluem todas as escritas/leitura de domínio, índices, histórico e tentativas; não reutilizar estimativa de ``1 leitura e 2 escritas'' como custo total.

\paragraph{Relatórios, transporte e integridade.}
O backend atual retorna PDF em base64; a especificação adota esse transporte para relatórios limitados, com conversão em Blob no cliente e liberação da URL temporária. Roteiros e comprovantes permanecem no Storage. Relatórios maiores exigem decisão de infraestrutura antes de ampliar limites. Um hash simples de parâmetros/data não é assinatura digital nem prova de autoria; tratá-lo como identificador de rastreabilidade. Relatórios históricos usam eventos/snapshots da época, não somente o cadastro atual.

\paragraph{Garantias de execução.}
Exemplos de código desta seção são ilustrativos e não substituem contratos UI/dicionário. Todas as leituras transacionais precedem escritas; callbacks podem repetir e não devem enviar e-mail ou gerar PDF como efeito externo. A transação grava evento/operação e um processador idempotente realiza efeitos externos após commit. Alterações no Auth e upload no Storage exigem reconciliação, pois não participam da transação Firestore. Consultas in e getAll não têm o mesmo contrato: tamanho de lote deve ser escolhido por API e medido; 200 não é limite universal do Firestore.

\paragraph{Ambiente e documentação operacional.}
README contém instruções de aplicação e emuladores; COMPILACAO\_NIX\_LCQUI.md é guia de build documental. O flake.nix atual fornece Node/Python/Git/Docker, mas não TeX Live; não presumir que nix develop disponibiliza pdflatex. Usar TeX Live instalado ou ambiente Nix dedicado e compilar main.tex até estabilizar referências. Os quatro documentos de acompanhamento registram evidências do código, não novos requisitos independentes.

\paragraph{Referências técnicas verificadas.}
\href{https://firebase.google.com/docs/firestore/manage-data/transactions}{Firebase: transações e lotes};
\href{https://firebase.google.com/docs/auth/admin/custom-claims}{Firebase: propagação de Custom Claims};
\href{https://firebase.google.com/docs/firestore/manage-data/enable-offline}{Firebase: persistência offline}.
Consulta em 11/09/2026. Estes fundamentos apoiam a execução transacional, renovação de tokens e opção por cache em dispositivo confiável.

\paragraph{Contratos dos auxiliares Q04/Q14.}
validarAceiteTcrTx lê o evento de aceite produzido em sessão autenticada do retirante, confere versão institucional vigente, UID, frasco, finalidade, operação ainda não consumida e justificativa de 20--2000 caracteres. Não aceita timestamp ou UID arbitrário no payload do gestor. vincularAceiteEAuditoriaTx consome o aceite e registra o empréstimo correspondente na mesma transação. validarAutoAtendimentoTx conta gestores ativos, bloqueia quando existe outro e exige justificativa; registrarNotificacaoChefiaTx grava notificação/outbox idempotente. resolverSnapshotQuimicoTx resolve unidade pelo resumo e densidade pela especificação, sem gravações antes de concluir leituras. Auxiliares ilustram contratos, não atestam implementação em functions/. O TCR não substitui segurança química institucional.
